\newif\ifuseseminar
\useseminartrue
\input{../../latex_common/header.tex.frag}

\title{Relatividade Restrita -- Lista de Exercícios: Dilatação do Tempo, Contração do Espaço, Paradoxo dos Gêmeos e Efeito Doppler}

\begin{document}

Em toda a lista abaixo usaremos assinaturas $(-,+,+,+)$ e as coordenadas em anos-luz.

\begin{enumerate}

      \item Tempo próprio e dilatação temporal na métrica de Minkowski: Considere duas
            trajetórias (curvas tipo-tempo) em um espaço-tempo plano com coordenadas
            $(ct, x^1, x^2, x^3)$:

            \begin{enumerate}
                  \item[Trajetória 1] partícula em repouso no referencial inercial:
                        $$x^1(t) = x^1_0,\quad x^2 = x^2_0,\quad x^3 = x^3_0.$$

                  \item[Trajetória 2] partícula em movimento uniforme na direção $x^1$:
                        $$x^1(t) = x^1_0 + v(t - t_0),\quad x^2 = x^2_0, \quad x^3 = x^3_0.$$
            \end{enumerate}

            \begin{enumerate}
                  \item Usando a métrica de Minkowski, calcule o tempo próprio
                        $\Delta \tau$ medido ao longo de cada trajetória no intervalo entre
                        $t_0$ e $t_1$.
                  \item Compare os dois tempos próprios e interprete o resultado como
                        uma dilatação temporal.
                  \item Mostre que o tempo próprio calculado ao longo da trajetória 2 é
                        igual ao tempo coordenado medido em um referencial inercial no qual a
                        partícula está em repouso. Esse resultado ilustra que o tempo
                        coordenado de um referencial é, por definição, o tempo próprio das
                        partículas em repouso nesse referencial.
                  \item Considere dois eventos $A$ e $B$. No referencial inercial $S$,
                        temos:
                        $$
                              x_A^\mu = (x^0_A, x^1_A) = (4, 1), \quad
                              x_B^\mu = (x^0_B, x^1_B) = (1, 5)
                        $$
                        Determine a velocidade $v$ de um referencial inercial $S'$ em
                        que os eventos $A$ e $B$ são simultâneos (i.e., ocorrem no mesmo
                        instante em $S'$).
            \end{enumerate}

      \item Paradoxo dos gêmeos e estrutura do espaço-tempo: Considere o paradoxo dos
            gêmeos, com um dos gêmeos partindo da Terra (posição $x = 0$) em uma viagem
            até $x = 2$ com velocidade $v = \sqrt{3/4}$, retornando com velocidade $v =
                  -\sqrt{3/4}$.

            \begin{enumerate}
                  \item Trace as linhas de simultaneidade nos referenciais inerciais da
                        Terra e da nave durante os trechos de ida e de volta. Mostre que
                        as linhas de simultaneidade da nave antes e depois da inversão
                        se cruzam. Discuta a consequência disso para a contagem do tempo
                        próprio e a interpretação dos eventos na perspectiva da nave.

                  \item A mudança de velocidade da nave implica uma descontinuidade nas
                        linhas de simultaneidade. Como descrever essa situação usando um
                        único sistema de coordenadas contínuo ao longo da trajetória da
                        nave?

                  \item Analise como a mudança instantânea de velocidade afeta o tempo
                        próprio acumulado pela nave. Como essa mudança pode ser
                        interpretada geometricamente no diagrama de espaço-tempo? E como
                        ela pode ser calculada?

                  \item Conceitualmente, por que não há simetria entre os dois gêmeos?
                        Qual elemento físico rompe a equivalência entre os referenciais?
                        Mostre como o uso de um único referencial inercial e o tempo
                        próprio ajuda a resolver o aparente paradoxo.
            \end{enumerate}

      \item Derivando o efeito Doppler relativístico: Considere dois observadores em
            movimento relativo com velocidade $v$ ao longo do eixo $x^1$.

            \begin{enumerate}
                  \item Defina o período de emissão de sinais do ponto de vista do
                        emissor e do observador. Explique como esse período pode ser calculado
                        no referencial de cada um, destacando a distinção entre tempo próprio
                        e tempo coordenado.

                  \item Suponha que o emissor está em movimento com velocidade $v$ e o
                        observador está em repouso à direita. Calcule a frequência
                        (própria) observada $\nu_o$ em função da frequência (própria)
                        emitida $\nu_e$, obtendo:
                        $$
                              \nu_o = \nu_e \sqrt{\frac{1 + \beta}{1 - \beta}}, \quad \beta = \frac{v}{c}
                        $$

                  \item Repita o cálculo no referencial em que o emissor está em repouso
                        e o observador se move com velocidade $-v$. Mostre que o resultado
                        obtido é o mesmo, reforçando a consistência do efeito Doppler
                        relativístico sob transformações de Lorentz.
            \end{enumerate}

      \item Vetores como derivadas direcionais:
            \begin{enumerate}
                  \item Mostre que a derivada direcional ao longo de um vetor $v$,
                        $$
                              \frac{\partial f}{\partial v} \equiv \lim_{\epsilon \to 0}
                              \frac{f(x + \epsilon v) - f(x)}{\epsilon},
                        $$
                        pode ser escrita na forma tensorial como:
                        $$
                              \frac{\partial f}{\partial v} = v^\mu \frac{\partial f}{\partial x^\mu}.
                        $$
                        Assim, o vetor $v$ pode ser representado por:
                        $$
                              v \equiv v^\mu \frac{\partial }{\partial x^\mu}.
                        $$

                  \item Dada uma transformação de coordenadas $y^\mu = \phi_y^\mu(x)$,
                        com inversa $x^\mu = \phi_x^\mu(y)$, use a representação anterior para
                        mostrar que as componentes transformadas do vetor satisfazem:
                        $$
                              v^{\prime\mu} = \frac{\partial \phi_y^\mu}{\partial x^\nu} v^\nu.
                        $$
                  \item Seja $x^\mu(t)$ uma curva parametrizada por $t$ em um espaço com
                        coordenadas $x^\mu$. Mostre que a variação infinitesimal de uma
                        função escalar $f(x)$ ao longo da curva pode ser escrita como
                        $$
                              \mathrm{d}f[x(t)] = \lim_{\epsilon \to 0} \frac{1}{\epsilon} \int_t^{t+\epsilon} \frac{\partial f}{\partial x^\mu} \frac{\mathrm{d}x^\mu(t)}{\mathrm{d}t} \,\mathrm{d}t = \frac{\mathrm{d}f}{\mathrm{d}v},
                        $$
                        onde $v^\mu \equiv \frac{\mathrm{d}x^\mu(t)}{\mathrm{d}t}$
                        representa o vetor tangente à curva. Conclua que a quantidade
                        $\frac{\partial f}{\partial x^\mu}$ define um objeto que age
                        linearmente sobre vetores: um \textbf{covetor} (ou 1-forma).

                  \item Com base na expressão anterior, mostre que a ação de um covetor
                        $w_\mu$ sobre um vetor $v^\mu$ pode ser escrita como:
                        $$
                              w(v) = w_\mu \,\mathrm{d}x^\mu(v) = w_\mu v^\mu.
                        $$

            \end{enumerate}
\end{enumerate}
\end{document}
